\pageof{Reading the rules}\label{sec:map}
Each circle is a laboratory that reports a classical outcome. We call a laboratory that sends an arrow a \emph{parent} of its recipient, and the recipient its \emph{child}. An arrow from one circle to another allows the recipient to use that outcome. Each diamond is a new, independent source. A quantum source can send different parts of one state to several laboratories. A classical source can hand each recipient a copy of a random string. Each laboratory also has private randomness.
\duo{Quantum source and two shares}{\vtx{a}{0}{0}{a}\vtx{b}{3}{0}{b}\qsrc{q}{1.5}{1.1}{a,b}}
{Classical source and two copies}{\vtx{a}{0}{0}{a}\vtx{b}{3}{0}{b}\csrc{q}{1.5}{1.1}{a,b}}
A graph has a \emph{quantum--classical gap} if some observed distribution can be produced using its permitted quantum resources and cannot be produced when every source is classical. A graph with an independently established gap is a \emph{seed}. We consider finite experiments and may choose the size of each laboratory's output. A laboratory can therefore report a pair, including a setting or a teleportation correction in addition to its usual result. A claim about a graph establishes the existence of a gap for some choice of output sizes.

\begin{rulebox}{Find a rule by the change you want to make}
\begin{tabular}{@{}p{.47\linewidth}p{.45\linewidth}@{}}
Fix a laboratory's announcement & Opening triangle example, p.~1\\
Upgrade a source or add connections & \readat{sec:resources}{Compare resources}\\
Forget an announcement & \readat{sec:forget}{Relay or teleport}\\
Keep a selected outcome & \readat{sec:select}{Postselection}\\
Identify a setting with an outcome & \readat{sec:stitch}{Stitch two nodes}\\
Give one node two recorded outputs & \readat{sec:stitch}{Split a node}\\
Remove inputs by perfect prediction & \readat{sec:template}{The shared argument}\\
Use a parent as predictor & \readat{sec:parent}{Parent Fritz}\\
Use separation in the original graph & \readat{sec:separate}{Original Fritz}\\
Combine independence statements & \readat{sec:kpc}{The KPC example}\\
Use an entropy inequality & \readat{sec:entropy}{The same KPC example}\\
Delete edges using independence alone & \readat{sec:nopred}{Prediction-free deletion}
\end{tabular}
\end{rulebox}
\textbf{Why these proofs can be chained.} Let \(P'\) be an arbitrary witness for the simpler graph. Construct a distribution in the original graph from which a specified operation recovers \(P'\). Prove that applying that operation to a classical explanation would explain \(P'\) classically. The contradiction proves a gap for the original graph. The next rule can now use that gap even if it constructs a very different witness.

\trio{Forget a relay}{\vtx{a}{0}{0}{a}\vn[pre]{v}{1.7,0}{$v$}\vtx{b}{3.4}{0}{b}\de{a}{v}\de{v}{b}
\draw[newedge] (a) to[bend right=40] (b);}
{Select a joint outcome}{\vtx{a}{0}{0}{a}\vtx{b}{3.4}{0}{b}\vn[pre]{v}{1.7,1.5}{$v$}
\qsrc{l}{0}{1.5}{a,v}\qsrc{r}{3.4}{1.5}{v,b}}
{Predict a field}{\vtx{x}{0}{0}{X}\vtx{s}{3.4}{0}{S}\csrc{l}{1.7}{1.5}{x,s}
\node[lab] at (1.7,-.65) {\(S=g(X)\)};}
\captext{Three ways to recognize the next move: follow an arrow through a relay, look for two sources meeting at a measurement, or look for common data that two laboratories can both announce.}

\pageof{When the same experiment still works}\label{sec:resources}
\textbf{Upgrade a source.} Replace a shared classical source by a quantum one on the same recipients. One can also give a subset of its recipients an additional quantum source. A gap already present survives both changes.
\duo{Triangle with an upgraded source}{\vtx{a}{0}{0}{A}\vtx{b}{3.8}{0}{B}\vtx{r}{1.9}{2.8}{R}
\qsrc{ab}{1.9}{-.7}{a,b}\qsrc{u}{-.15}{1.75}{a,r}\csrc{v}{3.95}{1.75}{b,r}}
{Known triangle seed}{\trianglefig{0}}
\quant A quantum source can draw a classical label and send perfectly distinguishable copies. An optional extra source can send a fixed state. The old experiment still works.
\clas After treating all sources as classical, the new graph gives exactly the old possibilities. In the subset case the larger source may distribute the extra independent string too, with the other recipients ignoring it.

\textbf{Add connections within a classical equivalence class.} A gap passes to a graph with extra arrows or larger sources if both graphs allow the same classical distributions. If both quantum and classical sets agree, a gap passes in either direction.
\duo{The graph with an extra arrow}{\bellsettings{1}}{Bell settings with a common cause}{\bellsettings{0}}
\quant Ignore the extra arrow. Every old experiment still runs.
\clas Here \(X\) takes input only from the settings source and its private randomness. That source can sample \(X\) in advance and give it to both settings laboratories. Thus it can simulate the added arrow in a classical model. Equality of the classical sets ensures the old witness is still outside them.
\textbf{When can a source work equally well in either form?} Suppose all but at most one recipient have no visible parent and receive no other shared source. Each such laboratory makes a fixed measurement. The source can make those measurements itself, send the outcomes as classical data, and tell the remaining recipient which conditional state to prepare locally. This preserves the quantum statistics when the source becomes classical. A free setting that reads a single source in the Bell example has precisely this property \cite{Equiv}.

\begin{ideabox}{Why check the classical sets?}
Add every forward arrow in a time ordering of the laboratories of an ordinary Bell experiment. The enlarged classical graph can sample \emph{any} observed distribution, one outcome at a time conditional on all its predecessors. Even a Bell violation has a classical explanation with those extra communications.
\end{ideabox}

\pageof{Forget an announcement}\label{sec:forget}
Forgetting \(v\) means summing over every value \(v\) reported. A child can still receive information that \(v\) forwarded during the experiment. In this example, 0 sends its announcement to 1 and shares a separate source with 2.
\trio{Before forgetting 0}{\margfig{0}}{Ordinary relay}{\margfig{1}}{With teleportation}{\margfig{2}}
\captext{The source between 1 and 3 stays quantum in both outputs. Look at the new link between 1 and 2: the ordinary relay supplies shared classical data. Teleportation also permits a quantum state there. The visible edge \(1\to2\) stays in all three pictures.}

\textbf{Ordinary marginalization.} Give each visible child of \(v\) all of \(v\)'s visible parents. Every incident source keeps its other recipients and its type. Add a classical source for its other recipients together with all of \(v\)'s children. Add shared classical randomness over the children for \(v\)'s private random choice. A classical source contained in a larger classical or quantum source can be absorbed by it.
\quant Run the desired resources for the surviving recipients. Each old source can also send a random string to \(v\) and its other recipients. Let \(v\) report the strings, its visible inputs and a private seed. Every child can extract what it needs.
\clas A child with these inputs can recompute the old announcement of \(v\). Sampling its noise once for all children preserves their joint distribution.

\textbf{Marginalization with teleportation.} Let \(T\) be the children of \(v\) that already share a quantum source with \(v\). Each quantum source incident to \(v\) may now also distribute quantum shares to every member of \(T\).
\quant Place each destination subsystem at \(v\). For a child in \(T\), its old shared quantum source can additionally supply an entangled pair. Laboratory \(v\) measures its half together with the subsystem and sends the classical correction through the visible arrow. The child recovers its subsystem. Distinct destinations receive distinct subsystems.
\clas The ordinary classical explanation just constructed still fits the enlarged sources. In the picture, making the 1--3 source classical in the rightmost graph reaches an instrumental seed: 3 supplies a setting to 1 through shared data, 1 sends its outcome to 2, and 1 and 2 share a quantum state.

\pageof{Keep only selected runs}\label{sec:select}
Choose a node \(v\) and an outcome of positive probability. Keep runs with that outcome, then drop \(v\). Let \(B\) be its visible parents and let its \emph{siblings} be the laboratories sharing a source with it. Add a classical source for \(B\) and all siblings together, and a quantum source for all its quantum siblings. Existing sources keep their surviving recipients.
\duo{Laboratory 0 measures two shares}{\swapfig{0}}{Conditioned on one Bell outcome}{\swapfig{1}}
\captext{Here 0 has no visible parents. A joint measurement at 0 can leave 2 and 3 entangled on the selected runs. The other source, between 1 and 2, survives.}
\begin{rulebox}{Two checks before postselecting}
Every visible grandparent of \(v\) must also be its parent. Every source touching \(B\) must either include \(v\) or have all its recipients inside \(B\).
\end{rulebox}
\quant One source at \(v\) draws the desired shared label; the others guess it uniformly. Keep agreements. Every label has the same chance of acceptance, so its desired law survives. If there is no such source, \(v\) draws the label. Prepare the new quantum state at \(v\), teleport the outside shares, and select the outcomes needing no correction.

The parents announce uniform guesses. Laboratory \(v\) can run their intended measurements in causal order: their visible parents lie in \(B\), sources wholly inside \(B\) can be prepared locally, and every source reaching both \(B\) and the outside can deliver the parents' shares to \(v\). Keep runs where the guesses match the measured outputs. For each possible output record the chance of a match is the same. The accepted records and outside systems thus have the desired joint statistics.
\clas Condition the old sources feeding \(v\) or lying inside \(B\) on the selected outcome. The new classical source can draw this joint sample and the conditional parent outputs together. Give each parent its coordinate and each sibling its original source data. Other sources remain independent, and children of \(v\) can insert its selected value.
\textbf{A failed check.} Each parent \(p_i\) and its outside observer \(o_i\) report the same fair bit, independently for the two pairs. Select when the parents agree. The outside observers now share a fair bit. In the proposed graph the new source reaches only the parents, so \(o_1,o_2\) must be independent. The original sources each violate the second check.

\duo{Select when the parents agree}{
\vtx{o}{0}{1.6}{o_1}\vtx{p}{1.4}{0}{p_1}\vtx{v}{2.8}{-1.1}{v}
\vtx{r}{4.2}{0}{p_2}\vtx{t}{5.6}{1.6}{o_2}
\csrc{l}{1.4}{1.6}{o,p}\csrc{k}{4.2}{1.6}{r,t}\de{p}{v}\de{r}{v}}
{The attempted replacement}{
\vtx{o}{0}{1.6}{o_1}\vtx{p}{1.4}{0}{p_1}
\vtx{r}{4.2}{0}{p_2}\vtx{t}{5.6}{1.6}{o_2}
\csrc{l}{1.4}{1.6}{o,p}\csrc{k}{4.2}{1.6}{r,t}\csrc{u}{2.8}{-1.1}{p,r}}

\pageof{Stitch two nodes or split one}\label{sec:stitch}
\textbf{Stitching.} Choose \(x\) with no visible parents and no shared source. Choose \(a\ne x\) outside its descendants. Delete \(x\) and let its children receive \(a\) instead. This condition keeps the new arrows acyclic. The search chooses a childless \(a\).
\duo{Bell experiment with free setting \(x\)}{\stitchfig{0}}{Use the value \(a\) as that setting}{\stitchfig{1}}
\quant In the Bell experiment sample \(x\) independently and uniformly from the alphabet of \(a\). Run the desired right-hand strategy, feeding \(x\) to the former children wherever it asks for \(a\). On the diagonal \(x=a\) every instruction receives precisely the intended value. For each outcome \(a\) and each record \(w\) at the remaining laboratories, the distribution in the stitched graph is
\[
   P'(a,w)=\frac{P(x=a,a,w)}{P_X(a)}=|\mathcal A|\,P(x=a,a,w),
\]
where \(\mathcal A\) is the set of possible values of \(a\). The denominator is constant because \(x\) was chosen uniformly. Each run in the stitched experiment is represented on the diagonal with the same weight.
\clas Fix all classical random choices except \(x\). The value of \(a\) is already determined independently of \(x\). Wire that value into the former children of \(x\). Averaging these runs gives exactly the same normalized diagonal. A gap on the right therefore proves one on the left.

\textbf{Splitting.} Replace \(s\) by \(s,s'\). Each inherits its visible parents, its outgoing arrows, and membership in each of its sources. Also give the pair shared classical randomness, which can be carried by any source they already share.
\duo{One output can be a pair}{\splitfig{0}}{Two visible coordinates}{\splitfig{1}}
\quant One laboratory can carry out both split laboratories' measurements on their respective systems and announce the pair. Each child can read either coordinate.
\clas With one laboratory reporting a pair, its conditional distribution given its inputs can be sampled from one private random seed. Give the split laboratories the same seed and inputs. They compute the same pair and each reports its own coordinate. Children again receive both. Thus a gap for the split graph proves a gap for the original. Fritz uses this when one coordinate records a setting and the other records a measurement result.
